% GENERATED FILE. Edit canonical lessons, then run npm run generate:books.
% canonical-source-hash: fnv1a64:e12b8a16ed25d3b2
% canonical-lessons: KA-C151-gellu, KA-C151-solu, KA-C151-baccidu, KA-C151-kapadu, KA-C151-sudu

\chapter{To win, To lose, To hide, To protect, To burn}
\label{ch:ka-to-win-to-lose-to-hide-to-protect-to-burn}
\hlchaptermodality{151}

\begin{chapteropening}
\textbf{By the end of this chapter:} \emph{I can say to win, to lose, to hide, to protect and to burn.}

Complete the last lesson of chapter 151: i can say to win, to lose, to hide, to protect and to burn.
\end{chapteropening}

\section[\texorpdfstring{\kn{ಗೆಲ್ಲು} (gellu)}{ಗೆಲ್ಲು (gellu)}]{\kn{ಗೆಲ್ಲು} (gellu) — to win}

\label{lesson:KA-C151-gellu}



\begin{quote}
Before the new one: say the Kannada for to tie, then the Kannada for to sew.
\end{quote}

\subsection*{You'll want to know: \kn{ಗೆಲ್ಲು}}
\textbf{\kn{ಗೆಲ್ಲು}} — \emph{gellu} — \textquotedblleft{}to win\textquotedblright{}.

\begin{grammarlens}[title={What you've built}]
One more word for this chapter.
\end{grammarlens}

\subsection*{Your turn}
\begin{itemize}
\raggedright
  \item \emph{Say it:} \emph{gellu}
  \item \emph{Say it:} \emph{gellu}, once more
  \item \emph{Say it:} say \emph{gellu}
  \item \emph{Recall:} say the Kannada for to tie, then the Kannada for to sew, then say \emph{gellu} again
\end{itemize}

\begin{tcolorbox}[breakable,colback=teal!4,colframe=teal!35!black,title={Before you move on}]
What does \kn{ಗೆಲ್ಲು} mean? (\textquotedblleft{}To win\textquotedblright{}.) Say it once more.
\end{tcolorbox}
\section[\texorpdfstring{\kn{ಸೋಲು} (sōlu)}{ಸೋಲು (sōlu)}]{\kn{ಸೋಲು} (sōlu) — to lose (a game)}

\label{lesson:KA-C151-solu}



\begin{quote}
Before the new one: say the Kannada for to sew, then the Kannada for to win.
\end{quote}

\subsection*{You'll want to know: \kn{ಸೋಲು}}
\textbf{\kn{ಸೋಲು}} — \emph{sōlu} — \textquotedblleft{}to lose (a game)\textquotedblright{}.

\begin{grammarlens}[title={What you've built}]
One more word for this chapter.
\end{grammarlens}

\subsection*{Your turn}
\begin{itemize}
\raggedright
  \item \emph{Say it:} \emph{sōlu}
  \item \emph{Say it:} \emph{sōlu}, once more
  \item \emph{Say it:} say \emph{sōlu}
  \item \emph{Recall:} say the Kannada for to sew, then the Kannada for to win, then say \emph{sōlu} again
\end{itemize}

\begin{tcolorbox}[breakable,colback=teal!4,colframe=teal!35!black,title={Before you move on}]
What does \kn{ಸೋಲು} mean? (\textquotedblleft{}To lose (a game)\textquotedblright{}.) Say it once more.
\end{tcolorbox}
\section[\texorpdfstring{\kn{ಬಚ್ಚಿಡು} (bacciḍu)}{ಬಚ್ಚಿಡು (bacciḍu)}]{\kn{ಬಚ್ಚಿಡು} (bacciḍu) — to hide}

\label{lesson:KA-C151-baccidu}



\begin{quote}
Before the new one: say the Kannada for to win, then the Kannada for to lose.
\end{quote}

\subsection*{You'll want to know: \kn{ಬಚ್ಚಿಡು}}
\textbf{\kn{ಬಚ್ಚಿಡು}} — \emph{bacciḍu} — \textquotedblleft{}to hide\textquotedblright{}.

\begin{grammarlens}[title={What you've built}]
One more word for this chapter.
\end{grammarlens}

\subsection*{Your turn}
\begin{itemize}
\raggedright
  \item \emph{Say it:} \emph{bacciḍu}
  \item \emph{Say it:} \emph{bacciḍu}, once more
  \item \emph{Say it:} say \emph{bacciḍu}
  \item \emph{Recall:} say the Kannada for to win, then the Kannada for to lose, then say \emph{bacciḍu} again
\end{itemize}

\begin{tcolorbox}[breakable,colback=teal!4,colframe=teal!35!black,title={Before you move on}]
What does \kn{ಬಚ್ಚಿಡು} mean? (\textquotedblleft{}To hide\textquotedblright{}.) Say it once more.
\end{tcolorbox}
\section[\texorpdfstring{\kn{ಕಾಪಾಡು} (kāpāḍu)}{ಕಾಪಾಡು (kāpāḍu)}]{\kn{ಕಾಪಾಡು} (kāpāḍu) — to protect, to save}

\label{lesson:KA-C151-kapadu}



\begin{quote}
Before the new one: say the Kannada for to lose, then the Kannada for to hide.
\end{quote}

\subsection*{You'll want to know: \kn{ಕಾಪಾಡು}}
\textbf{\kn{ಕಾಪಾಡು}} — \emph{kāpāḍu} — \textquotedblleft{}to protect, to save\textquotedblright{}.

\begin{grammarlens}[title={What you've built}]
One more word for this chapter.
\end{grammarlens}

\subsection*{Your turn}
\begin{itemize}
\raggedright
  \item \emph{Say it:} \emph{kāpāḍu}
  \item \emph{Say it:} \emph{kāpāḍu}, once more
  \item \emph{Say it:} say \emph{kāpāḍu}
  \item \emph{Recall:} say the Kannada for to lose, then the Kannada for to hide, then say \emph{kāpāḍu} again
\end{itemize}

\begin{tcolorbox}[breakable,colback=teal!4,colframe=teal!35!black,title={Before you move on}]
What does \kn{ಕಾಪಾಡು} mean? (\textquotedblleft{}To protect, to save\textquotedblright{}.) Say it once more.
\end{tcolorbox}
\section[\texorpdfstring{\kn{ಸುಡು} (suḍu)}{ಸುಡು (suḍu)}]{\kn{ಸುಡು} (suḍu) — to burn}

\label{lesson:KA-C151-sudu}



\begin{quote}
Before the new one: say the Kannada for to hide, then the Kannada for to protect.
\end{quote}

\subsection*{You'll want to know: \kn{ಸುಡು}}
\textbf{\kn{ಸುಡು}} — \emph{suḍu} — \textquotedblleft{}to burn\textquotedblright{}.

\begin{grammarlens}[title={What you've built}]
One more word for this chapter.
\end{grammarlens}

\subsection*{Your turn}
\begin{itemize}
\raggedright
  \item \emph{Say it:} \emph{suḍu}
  \item \emph{Say it:} \emph{suḍu}, once more
  \item \emph{Say it:} say \emph{suḍu}
  \item \emph{Recall:} say the Kannada for to hide, then the Kannada for to protect, then say \emph{suḍu} again
\end{itemize}

\begin{tcolorbox}[breakable,colback=teal!4,colframe=teal!35!black,title={Before you move on}]
What does \kn{ಸುಡು} mean? (\textquotedblleft{}To burn\textquotedblright{}.) Say it once more.
\end{tcolorbox}
